\clearpage
\fancyhead[L]{\small D'AGUILA.}
\hypertarget{native-546-middle}{}\label{native-546-middle}
\begin{center}\small Göttingische gelehrte Anzeigen. No. 102. pp. 1009..1010. June 27, 1812.\\[3pt]\rule{.72\linewidth}{.3pt}\end{center}

\emph{Decouverte de l'orbite de la terre du point central de l'orbite du soleil, leur situation et leur forme; de la section du zodiaque, par le plan de l'équateur, et du mouvement concordant des deux globes. Avec figures. Par Mr.} \textsc{C. J. E. H. d'Aguila}, \emph{ancien élève du génie.} 1806. 432 pages in octavo. Paris. \emph{De l'imprimerie de Délance.}

As requested, we now give a brief review of this book. It appeared 6 years ago but only recently came to our attention. Here we have a new anti-Copernican who has the Sun describe a circle around the Earth every year and the planets move around the Sun (although he mentions the planets only briefly and in passing). In addition to the Earth's rotation, however, he assigns it a uniform annual motion in an eccentric circle in the plane of the equator. He sets the radius of this circle at $\frac{1}{59}$ of the radius of the Sun's annual orbit. This important discovery, which the author himself calls incomprehensible and attributes to direct inspiration (which is also why he dedicates it to God himself), is to glorify the nineteenth century and destroy the sects of \textsc{Copernicus}, \textsc{Kepler} and \textsc{Newton}, on which the author looks down with deep contempt. That may be enough to recommend the book to inquisitive readers, if only as a psychological curiosity.
